% --------------------------------------------
% CHAPTER 2: THEORETICAL BACKGROUND
% --------------------------------------------

\chapter{Theoretical Background}\label{chap:background}

Due to the complicated nature of this work, it is necessary to provide a theoretical background for some topics fundamental to the research. While not completely necessary to appreciate the value of this work, readers are recommended to use this section as a reference for confusing topics. The organization of this chapter builds from topics on crystallography to \gls{qm} fundamentals, to the more complicated underlying theories that define the software used in this work. % Additional readings are provided in the appendices. symbols can be found in the glossary. ???

% ============================================================
\section{Crystallography and Periodic Systems}
% ============================================================

Crystallography is the experimental science of the arrangement of atoms in solids. Given the appropriate conditions, groups of atoms crystallize into a repeated arrangement with an observable structure and set of properties. A few crystallography concepts are needed to support the more complex topics discussed later. The \textit{crystal lattice} is a lattice of atoms in a crystal that have a pattern that repeats periodically. The \textit{unit cell} is the smallest repeating structural unit within a crystal lattice that shows the full symmetry of the crystal structure \autocite{westSolidStateChemistry2022}. The repeating nature of a periodic system is key for atomistic codes, since it allows for a smaller simulation box with assumed periodic boundary conditions, rather than having to directly simulate thousands of atoms. 

In this work, we assume that the materials of interest are inorganic solid crystals, though in reality they could exist in the liquid or gaseous phases. Thanks to this assumption, we can enforce periodic boundary conditions onto the simulation. This means that a unit cell or a few unit cells together can be used to represent the crystal as a whole, since the lattice can be expected to repeat along vectors known as the crystallographic basis. The primitive lattice, or Bravais lattice, is defined as an infinite array of discrete points generated by the crystallographic basis vectors which span the lattice \autocite{ashcroftSolidStatePhysics1976}. There are 14 unique Bravais lattices that are used to define crystalline materials. The primitive cell is defined as the volume of space that exactly fills all the space in the lattice when translated through all the vectors in a Bravais lattice.

\subsection{Reciprocal Space and the Brillouin Zone}

The reciprocal lattice is a crucial concept for simulation, that leads to many interesting conclusions. The Solid State Physics textbook from Ashcroft and Mermin has a good introduction to reciprocal space. 
"Consider a set of points constituting a Bravais lattice \textbf{R}, and a plane wave $e^{ik \cdot r}$. For general k, such a plane wave will not, of course, have the periodicity of the Bravais lattice, but for certain special choices of wave vector it will. The set of all wave vectors \textbf{K} that yield plane waves with the periodicity of a given Bravais lattice is known as its reciprocal lattice"\autocite{ashcroftSolidStatePhysics1976}. So the set of wave vectors \textbf{K} for all points \textbf{R} will satisfy $$e^{iK \cdot  R}=1 $$
Reciprocal lattices are defined with reference to a particular Bravais lattice.

 The Wigner-Seitz primitive cell of the reciprocal lattice is called the first Brillouin zone (BZ). Many of the calculations in this work involve integrating periodic functions of a Block wave vector over specified portions of the BZ. Monkhorst-Pack developed a method for finding "special" points to integrate over within the BZ in order to optimize calculations \autocite{monkhorstSpecialPointsBrillouinzone1976}. These "special" points are known as k-points. A high-symmetry path is then defined between these k-points, called the k-path. 

% ============================================================
\section{Quantum Mechanics Basics}
% ============================================================

Quantum mechanics is a famously difficult subject, so this section will attempt to introduce the bare minimum concepts needed to support the Chapter \ref{chap:methods} methods overview. 

We start with electrons, fundamental subatomic particles with a small mass and a negative electrical charge equal to $1.062 \times 10^{-19}$ Coulomb. The number of electrons in an atom is typically equal to the number of protons (positive subatomic particles that determine the type of atom), but when this number is not equal, the atom is called an ion. Electrons occupy electronic orbitals outside the nucleus of an atom. These orbitals define regions in which electrons have the highest probability of being found. This concept of probability in relation to subatomic particles will be discussed later. There are two general categories of electrons, core electrons that occupy the lowest energy levels, and valence electrons, which occupy the outermost shell and have the highest energy levels. Valence electrons primarily determine the reactivity of the atom and engage in chemical bonding. When all of a system's electrons occupy the lowest available energy levels, the atom is said to be in its ground state. The ground state is especially important in chemistry, since it is used to predict the crystal's reactivity, binding capabilities, and other properties. When the atom is not it's ground state, it is said to be in an excited state.

Phonons are quasi-particle representations of vibrational energy in a solid material. Though not physical particles in reality, they can collide and interact with other particles, and play a large role in how kinetic energy moves through a crystal lattice. Phonons are divided into two modes, acoustic and optical modes. Acoustic mode phonons are in phase, meaning that they move together in the same direction. They vanish at the (0,0,0) k-point, and correspond to sound waves propagating through the lattice. Optical mode phonons are out of phase, meaning that positive and negative ions move in opposite directions. The optical phonons are related to electromagnetic radiation. Both types contribute to heat transport an electron scattering within the lattice.
In order to properly discuss phonons, we must mention the harmonic approximation, and the effects of anharmonicity.
Thermal properties such as heat capacity, entropy, and thermal conductivity can be predicted using phonon calculations.

Define phonons, and introduce their connection to thermal conductivity (as a mention not explanation)
quantized vibrations in the form of phonons
Harmonic, anharmonic

\autocite{togoFirstprinciplesPhononCalculations2023}
% → Extract: Supercell finite-displacement method for 
%   IFC extraction from DFT forces. This is the 
%   theoretical basis for your Phonopy workflow.
\autocite{togoImplementationStrategiesPhonon2023}
% → Extract: Symmetry reduction of IFC calculations; 
%   cite for completeness alongside above.

% What is a band structure? What makes a band structure diagram useful? These are crucial outputs for my dft calculations. Mention fermi energy, band gap...
% How is band structure used to quantify and display electron and phonon results?
% What is the density of states? How is it useful? How is it related to the 3D density of states generated by Dr. Alex

% ============================================================
\section{Thermodynamic Formulation}
% ============================================================

At the center of this work is the attempt to predict the thermal properties of inorganic solids using atomic scale calculations. Below is a brief overview of the mathematical foundation for deriving these properties.

\subsection{Fourier's Law}

According to the 0th law of thermodynamics, any temperature gradient ($\nabla T$)\glsadd{gradient}\glsadd{temperature} will spontaneously induce a heat flux \gls{heat flux} to return the system to thermodynamic equilibrium.

Fourier's Law describes the relationship between the heat flux \gls{heat flux} and the temperature- and pressure-dependent thermal conductivity \gls{kappa}$(T,p)$\glsadd{temperature}\glsadd{pressure}

\begin{equation}
    \boldsymbol{J} = -\boldsymbol{\kappa}(T,p)\cdot \nabla T
\end{equation}

\subsection{Heat Capacity}

The isochoric specific heat capacity (\gls{specific heat capacity}) is defined as the amount of heat that must be supplied to a mass unit of a substance at constant volume to raise its temperature by some amount. In general, the heat capacity of a substance under a temperature differential ($dT$) with some mass ($m$) and supplied heat ($dQ$) can be shown as:
\begin{equation}
    c_v=\frac{dQ}{m\,dT}
\end{equation}
In the more specific case of phonon simulation, specific heat capacity can be defined in terms of the phonon frequency ($\omega$)

\autocite{ashcroftSolidStatePhysics1976} % ch 23 phonon contribution to heat capacity (this is where I pulled my eq)

\subsection{Green-Kubo Theory}

The paper Carbogno et al \autocite{carbognoInitioGreenKuboApproach2017} presents a method for the first-principles formulation of the Green-Kubo method for phonon thermal conductivity of solids in equilibrium. 

The Green-Kubo formula for heat flux developed for LAMMPS using the stress (\gls{stress}) of a system of $i$ atoms \autocite{surblysApplicationAtomicStress2019}.

\begin{equation}
    \boldsymbol{J}=\frac{1}{V}\left[ \sum_i E_i \cdot \vec{v}_i - \sum_{i}\boldsymbol{S} \cdot\vec{v}_i \right]
\end{equation}
Which sums the convective flux $J_\text{conv}=E_i \cdot \vec{v}_i$ and the conductive flux $J_\text{cond}=-\boldsymbol{S} \cdot\vec{v}_i$.

%Still need to discuss heat flux autocorrelation function (HFACF) from Carbogno

\autocite{knoopInitioGreenKuboSimulations2023} % green kubo simulation in action, this could also go in lit review

\autocite{surblysApplicationAtomicStress2019} % this is the LAMMPS heat-flux equation paper, gives a look into atomic stress formulation
\autocite{surblysMethodologyMeaningComputing2021} % the other LAMMPS heat-flux paper that has the other stress form

\autocite{zhangGreenKuboFormalismThermal2023} % this is part of the future work but might be useful to show for completeness

\subsection{Boltzmann Transport Equation}

What is BTE, and how does it connect to smaller scale mathematical formulas?
\autocite{ashcroftSolidStatePhysics1976} % ch 13 BTE for phonons, thermal conductivity
\autocite{harterPredictingMesoscaleSpectral2020} % this is basically the foundational paper for my results, need to pull from their derivations
Maybe use the Rattlesnake formulas here? I don't want to get that deep into it.
\autocite{harterCharacterizationThermalConductivity2015} % another jackson paper if I need it

\clearpage